\cotefile{157-5}
\scanpage{603}\mspage{12}
% archive sheet number: 602

\[
\underset{\substack{\Vert\\ \mathrm{Ens}\smallsetminus\{\emptyset\}}}{\mathrm{Ens}^{*}}
 \subset \mathrm{Ens}
\]
des ensembles non vides, on sait qu'elle est \unsure{accessible},
\deleted{\unsure{par conséquent}} mais elle n'est pas stable par
$\varprojlim$ de type $L_{\pi}$, \unsure{quelque|quelques} \unsure{grand} qu'on
\unsure{prenne} $\pi$.

Considérons
\[
\underbrace{\underline S=\underline S_a}_{\Vert\,\text{\unsure{déf|def}}}\subset M
\]
\[
f^{-1}(\mathrm{Ens}^{*}),
\]
sous-catégorie str. pleine de $M$ formée des $x\in\mathrm{Ob}\,M$ tels que
\[
a(x)\ne\emptyset \qquad\text{i.e.}\qquad x\leq a:\quad
\text{pour le préordre can. de }M.
\]

C'est donc une sous-catégorie \uline{ouverte} de $M$.
\deleted{\unsure{S}}Son complémentaire \unsure{fermé|qui en}
\inserted{$\underline S'=\underline S'_a$} est la ss/catégorie
fermée image inverse, par $f$, de la sous-catégorie str. pleine
\deleted{$\unsure{\{\emptyset\}}$} de $\mathrm{Ens}$ réduite au seul objet
$\{\emptyset\}$. Cette sous-catégorie de $\mathrm{Ens}$ est
\deleted{\unsure{une catégorie}} \inserted{la catégorie} finale type.
Elle est stable dans $\mathrm{Ens}$ par $\varinjlim$ quelconques, et
par $\varprojlim$ filtrantes quelconques. \deleted{\unsure{A}}Ainsi,
$\underline S'$ est stable \inserted{dans $M$} par $\varinjlim$
filtrantes quelconques, en particulier elle est $L_{\pi}$-stable pour
$\pi$ \unsure{n'imp.} quel. Explicitons :
\scanpage{604}\mspage{13}
% archive sheet number: 603

\[
\underbrace{\underline S'=S'_{\unsure{a|\cdot}}}_{\Vert\,\text{\unsure{def|déf}}}\subset M
\]
\[
\begin{aligned}
 f^{-1}(\{\emptyset\})={}&\text{ss-cat. \deleted{\unsure{str. pleine de|fermée de}} }M\\
 &\text{formée des }x\in\mathrm{Ob}\,M\text{ tels}\\
 &\text{que } a(x)=\emptyset .
\end{aligned}
\]

Comme $\underline S'$, \deleted{\unsure{étant}} fermée dans $M$, est
stable par $\unsure{\varinjlim|\underline{\lim}}$ non vides et en
particulier $L_{\pi}$-stable pour $\pi$-\unsure{gros|gd}, il y a plus de raisons de
s'attendre à ce que $\underline S'$ soit accessible dans $M$ que pour
\inserted{son complémentaire} $\underline S$, \deleted{\unsure{qui}}
dont il n'est pas clair qu'il ait cette propriété de stabilité, mais
seulement stabilité par $\varprojlim$ non vides, et en particulier par
facteurs directs.

En fait, nous allons voir qu'il peut arriver que $\underline S$ et
$\underline S'$ soient tous deux non accessibles dans $M$ :

\uline{Proposition} \uline{4} \unsure{Soient|Soit} $\mathcal U$ un univers, contenant
un cardinal infini $\pi_0$, tel que pour tt $E\in\mathcal U$,
\deleted{\unsure{tt}} \deleted{\unsure{tout}} $\pi_0$-ultrafiltre sur
$E$ soit trivial. Soit alors $k\in\mathcal U$ un corps de cardinal
\inserted{$\pi'_0$}$\geq\pi_0$, et soit $M$ la catégorie
\scanpage{605}\mspage{14}
% archive sheet number: 604

des $k$-algèbres commutatives. Donc $M$ est un
$\mathcal U$-\unsure{paratopos|paratope}, admettant $k$ comme objet initial.
Soit $\underline S'$ la sous-catégorie \unsure{strictmt} pleine,
stable par \inserted{petites} $\unsure{\varinjlim|\underline{\lim}}$
\deleted{quelconques}\inserted{non vides}, formée des objets $A$ de
$M$ tels que $\mathrm{Hom}_M(A,k)=\emptyset$, et $\underline S$ la
\deleted{\unsure{s}}sous-catégorie str. pleine complémentaire, formée des objets $A$ tels
que $\mathrm{Hom}_M(\deleted{\unsure{A}}A,k)=\emptyset$. Alors $\underline S'$ est une
\deleted{\unsure{fer}}sous-catégorie \uline{fermée} de $M$, et
$\underline S$ la sous-catégorie ouverte complémentaire,
\inserted{donc $\underline S$ (tout comme $\underline S'$) est}
stable par facteurs directs. Ceci posé,

a) \deleted{\unsure{Pour tout}} \inserted{Il n'existe pas de}
cardinal $\pi$ tel que $\underline S$ soit $L_{\pi}$-stable dans
$M$, a fortiori $\underline S$ n'est pas accessible dans $M$.

b) Bien que $\underline S'$ soit $L_0$-stable dans $M$
% slanted
\margin{Pas \unsure{prouvé} ?}
(\inserted{\unsure{ie}} stable par $\unsure{\varinjlim|\underline{\lim}}$
filtrantes, et \unsure{$=$}, on l'a dit, par ttes $\varinjlim$ non
\unsure{vide|vides}), $\underline S'$ n'est pourtant pas accessible dans $M$.

\uline{Démonstration} \deleted{Soit $E$ un} L'assertion a) va résulter du

\uline{Corollaire 1} \deleted{\unsure{Soit}}Sous les conditions de prop. 4, soient
$E\in\mathcal U$ un ensemble \inserted{infini}, $\Phi_0$ le filtre des
complémentaires des parties finies de $E$. Alors
\deleted{$\unsure{k_{\Phi}}$} $\mathrm{Hom}_M(k_{\Phi_0},k)=\emptyset$,
a fortiori pour tt filtre $\Phi$ sur $E$ plus fin que
$\unsure{\Phi|\Phi_0}$, \unsure{on a} $\mathrm{Hom}_M(k_{\Phi},k)=\emptyset$.
\scanpage{606}\mspage{15}
% archive sheet number: 605

On a posé ici, bien sûr, pour tt filtre $\Phi$ sur $E$
\[
 k_{\Phi}=\varinjlim_{J\in\Phi}k^J .
\]

On \deleted{\unsure{s}}sait que \inserted{pour tte $k$-algèbre comm. $A$,}
$\mathrm{Hom}_M(\deleted{\unsure{k_{\Phi}}}A_{\unsure{\Phi|\deleted{\Phi}}},k)$ est en
corr. 1-1 avec la partie de $\mathrm{Max}(\deleted{\unsure{k}}A_{\Phi})$
formée des idéaux maximaux $m$ de $\deleted{\unsure{k}}A_{\Phi}$
tels que \deleted{$\unsure{A}$} $k\unsure{-|\to}A/m$ soit un iso.
Si $A\inserted{\text{ est de la forme }}\unsure{=}k_{\Phi}$, $\Phi$ un filtre
sur $E$, cet ensemble est en corr. bi-univoque \deleted{\unsure{a}}avec l'ens des
\inserted{$\pi'_0$-}ultrafiltres $\Psi$ sur $E$ plus fins que $\Phi$,
où $\pi'_0=\mathrm{Card}(k)$ (proposition 3, et cor. 3
\inserted{de prop. 2} page 7). \deleted{\unsure{Comme}} vu
$\pi'_0\geq\pi_0$, un tel ultrafiltre est trivial
\inserted{par hyp.}.

Donc on trouve
\[
\begin{aligned}
\mathrm{Hom}_M(k_{\Phi},k)&\simeq\text{ensemble des }\xi\in E\text{ tels que}\\
&\{\xi\}\in\deleted{\Phi}\bigcap_{A\in\Phi}A
\end{aligned}\tag*{(*)}
\]
Mais \deleted{\unsure{si}} \inserted{dire que} $\Phi$ est plus fin que
$\Phi_0$ signifie aussi que pour tt $\xi\in E$, on a
$\deleted{\unsure{A}}E\smallsetminus\{\xi\}\in\Phi$, i.e.
$\exists A\in\Phi$, \unsure{avec|où} $A\subset E-\{\xi\}$ i.e.
\deleted{ce qui implique $\{\xi\}\notin\Phi$ (}
\deleted{\unsure{donc}}\inserted{i.e.} $\{\xi\}\notin\bigcap_{A\in\Phi}A$.
Donc on a alors $\mathrm{Hom}_M(k_{\Phi},k)=\emptyset$.

\deleted{Notons maintenant que} La conclusion du cor. 1 s'exprime
par la relation
\[
 k_{\Phi}\notin\mathrm{Ob}\,\underline S\qquad\text{i.e.}\qquad
 k_{\Phi}\in\mathrm{Ob}\,\underline S',
\]
d'autre part on a par définition
\[
 k_{\Phi}=\varinjlim_{J\in\Phi}k_{\unsure{J|D}},
\]
\scanpage{607}\mspage{16}
% archive sheet number: 606

et \deleted{\unsure{on}} il est clair (comme les $J\in\Phi$ sont
$\ne\emptyset$) que \deleted{pour} les $k_J$ en question sont dans
$\underline S$. Ainsi, pour un filtre $\Phi$ comme dans le cor. 1,
on trouve un système inductif d'objets de $\underline S$, indexé
par $\Phi^{\unsure{o|\circ}}$, dont la limite inductive n'est pas dans
$\underline S$. Donc le a) de la proposition 4 se déduit en notant
que pour tt cardinal $\pi$\inserted{$\in\mathcal U$}, on peut trouver
un ens $E\in\mathcal U$, et un filtre \inserted{$\Phi$} sur $E$ plus
fin que $\Phi_0$, tel que l'ens ordonné $\Phi^o$ soit \unsure{cord|croit|grand}
devant $\pi$. Pour \deleted{\unsure{t}}trouver un tel $\Phi$, il suffit de prendre $E$ de
cardinal $\pi'>\pi$, et pour $\Phi$ l'ens des complémentaires des
parties de $E$ de cardinal $\leq\pi$, de sorte que
$\unsure{\Phi^o|\Phi_0}$ est isom. à l'ens ordonné des
\deleted{\unsure{classes}} parties \inserted{de $E$} de cardinal
$\leq\unsure{\pi|\chi}$, lequel est bien \unsure{gros} devant $\pi$.
Cela prouve \deleted{\unsure{facilt}} le a) de prop. 4.

Prouvons maintenant le b). \unsure{Admettons} le

\uline{Lemme} Soit $M$ une catégorie accessible stable par
$\varinjlim$ \deleted{filtrantes dénombrables}
\inserted{de type $I_\alpha$, $\alpha$ un ordinal limite donné},
$\underline S'$ une sous-catégorie \uline{fermée}, $\pi$ un cardinal utile
\scanpage{608}\mspage{17}
% archive sheet number: 607

pour $M$, \inserted{tel que $\mathrm{card}\,\alpha\leq\pi$.}
Pour \deleted{que $\underline S'$ soit} tout
$x\in\mathrm{Ob}\,\deleted{\unsure{M}}{\underline S'}$, choisissons
une représentation \deleted{$\unsure{\alpha}$}-$\pi$-adaptée
($\omega$ étant le premier ordinal infini)
\[
 x\simeq\varinjlim_{i\in I_x}x_i
\]
(avec $I_x$ \unsure{gros} devant $\pi$, stable par
$\unsure{\varinjlim|\underline{\lim}}$ \deleted{filtr. dénombrables}
\inserted{de type $I_\alpha$}, $i\longmapsto x_i:I\unsure{-|\to}M$
commute aux dites limites, enfin les $x_i\in\mathrm{Ob}\,M(\pi)$).
Pour que $\underline S'$ soit $\pi$-accessible dans $M$, il f. et s.
que pour tt $x\in\mathrm{Ob}\,\underline S'$, il existe un
$i\in I_x$ tel que $x_i\in\mathrm{Ob}\,\underline S'$.

En effet, \deleted{c'est néces} notons que $\underline S'$
\deleted{étant} fermée, est stable \inserted{dans $M$} par tte
$\varinjlim$ non vide existt dans $M$, et en particulier par
$\varinjlim$ \deleted{filtrantes} \inserted{de type $I_\alpha$}.
Donc on peut appliquer T.13.22 \inserted{(p.342)} à
$M,\underline S',\pi$ et $\Phi=\{I_{\omega}\}$.
Alors $\unsure{a|d})\Longrightarrow c)$ implique que si
$\underline S'$ est $\pi$-accessible dans $M$, alors pour tt
$x\in\mathrm{Ob}\,\underline S'$, \unsure{alors l'ens|l'ens} $I_{x,S'}$ des $i\in I_x$
tels que $x_i\in\mathrm{Ob}\,\underline S'$ est cofinal dans
$\deleted{\unsure{J}}I_x$, \deleted{\unsure{or}} Cela implique qu'il
est non vide. (Et réciproquement d'ailleurs $\underline S'$ étant un
fermé \unsure{dans|de} $\deleted{\unsure{S}}M$, puisque si $x_i\in\underline S'$,
alors pour $j\unsure{\subset|\in}\mathrm{Ob}\,I_x$ avec $i\leq j$,
on a $\unsure{x_j|y}\in\underline S'$ \deleted{\unsure{donc}}
\inserted{i.e.} $I_{x,S'}$ contient l'image can. de $i\backslash I_x$,
laquelle est cofinale dans $I_{\unsure{\ldotp\ldotp\ldotp|x}}$).
Cela prouve le ``il faut''.
Prouvons le ``il suffit''. \unsure{Donc}
prouvons que $\underline S'$ est égal à l'image essentielle de
$\mathrm{Ind}_{\pi}(\underline S'(\pi))$. Elle la contient,
puisque $\underline S'$ est stable par $\varinjlim$ de type $L_{\pi}$.
\scanpage{609}\mspage{18}
% archive sheet number: 608

Donc il reste à prouver que tt $x\in\mathrm{Ob}\,\underline S^{\unsure{*|\prime}}$
est dans l'image essentielle. Or l'hypothèse nous
\deleted{\unsure{dit}}\inserted{assure} que
$x\simeq\varinjlim_{i\in\deleted{\unsure{\Phi}}I_x}x_i$ peut s'écrire
$\unsure{\varinjlim|\underline{\lim}}_{i\in I_{x,\underline S'}}x_i$
(car $I_{x,\underline S'}$ est cofinal dans $I_x$), où cette fois les
$x_i$ sont dans $\underline S'$. Comme $I_x$ \unsure{gros} devant
$\pi$, de $=$ la sous-catégorie pleine cofinale $I_{x\underline S}$,
donc on \unsure{gagne|y gagne}.

\uline{Corollaire} \deleted{\illegible{1}} Sous les conditions
précédentes, si $\underline S'$ est $\pi$-accessible, il en est de
$=$ de l'ouvert complémentaire $\underline S$, \uline{pourvu} que
$\underline S$ soit stable \inserted{dans $M$} par $\varinjlim$ de type
$L_\pi$\margin{!}

Il faut prouver que $\underline S$ est égale à l'image essentielle
de $\mathrm{Ind}_{\pi}\,\underline S(\pi)$.
\deleted{Soit donc}\inserted{L'hypothèse ``\unsure{pourvu}''}
\deleted{Prouvons}\inserted{assure} \deleted{d'abord}
que cette image essentielle est $\subset\underline S$.
\deleted{\unsure{Inversement}, soit donc
$\left|\,x=\varinjlim_I x_i\quad
\left|\begin{aligned}&x_i\in\mathrm{Ob}\,\underline S(\pi)\\
&I\ \text{\unsure{gros} devant }\pi\end{aligned}\right.\right)$,}
prouvons que \deleted{prouvons que $x\in\mathrm{Ob}\,\underline S$.
Sinon, on aurait} \deleted{\unsure{pour}}
$x\in\mathrm{Ob}\,\underline S$ est dans cette image essentielle.
On aura
\[
 x\simeq\varinjlim_{i\in I}x_i\qquad
 \left|\begin{aligned}&x_i\in\mathrm{Ob}\,M_\pi\\
 &I=M_\pi/x\\&\text{\unsure{gros} devant }\pi\end{aligned}\right.,
\]
et comme $x\in\mathrm{Ob}\,\underline S$ et $\underline S$ ouvert,
les $x_i$ \deleted{prouvons que $\varinjlim x_i\in$} sont dans
$\underline S$, donc dans $\underline S(\pi)$, OK.

\deleted{\uline{\unsure{Remarque}}} Dans le cas qui nous occupe,
\inserted{où $M=(k\text{-alg})$,} comme $\underline S$ n'est
justement pas stable dans $M$ pour $L_\pi$
\inserted{(quel que soit $\pi$)}, \unsure{(c'est|c'est} \unsure{cela qui|cela que} nous faisait conclure que
\scanpage{610}\mspage{19}
% archive sheet number: 609

$\underline S$ n'est pas accessible dans $M$), on ne peut donc
conclure de cette inaccessibilité de $\underline S$, celle de
$\underline S'$. Notons le

\uline{Corollaire} Soit $\pi$ un cardinal infini. Pour que la
sous-catégorie fermée $\underline S'$ de
$\deleted{\unsure{\underline S}}\unsure{M|\underline M}$ (formée
des \deleted{$A\in\mathrm{Ob}\,M$} $k$-alg. comm. $A$ telles que
$\mathrm{Hom}_M(A,k)=\emptyset$) soit $\pi$-accessible, il f et s
que pour tte $k$-algèbre $A$ [de cardinal $>\pi$],
\inserted{telle que $\mathrm{Hom}_M(A,k)=\emptyset$,}
\unsure{existe} une sous-algèbre $A_i$
\uline{de cardinal} $\leq\pi$, telle que
$\mathrm{Hom}_M(A_i,k)=\emptyset$.

(On applique le lemme, en prenant pour $\alpha$ \unsure{le} premier
ordinal infini.)

\uline{Question} Cela est-il vérifié, en prenant \inserted{\unsure{simplement|simplemt}}
$\pi\geq\mathrm{Card}\,k$ ? La question de l'existence
d'un cardinal $\pi$, ayant la propriété indiquée plus haut me paraît
délicate, et j'ignore la réponse. Si on veut construire un
contre exemple pour tt $\pi$, les développements précédents suggèrent
de prendre pour $A$ un produit infini d'extensions non triviales de
$k$, ou un \deleted{\unsure{corps}}\inserted{\unsure{anneau}} $k_{\Phi}$ comme dans
la dém. de prop. 4 a).
\scanpage{611}
% archive sheet number: 610

\uline{$\overline{\mathrm{XIX}}$}

Modèles\unsure{.} (4)
\scanpage{612}
% archive sheet number: 611

\diagram{
\(\begin{array}{@{}r l r@{}}
1 & \text{Rappels sur les couples de \unsure{Quillen}} & \text{1-12}\\
2. & \text{Rappels sur les triples de Quillen. Triples \uline{spéciaux}}
   & \text{13-32}\\
 & \text{(th. 2.3 p.19, th. 2.6 p. \unsure{25|21}, th. \unsure{2.9|2.7} \unsure{p.ge|page} 28)} &\\
3. & \text{Catégories à cofibrations} & \text{\deleted{\unsure{3}}\unsure{33|35}-50}
\end{array}\)
}
\scanpage{613}\mspage{1}
% archive sheet number: 612

1. \uline{Rappels sur les couples de Quillen}. \uline{Couples de Quillen spéciaux}.

Je vais reprendre les notations et la terminologie pour les relations de
LLP et RLP de Quillen. J'écris (pour une catégorie $M$, et
$\Phi,\Psi\subset\mathrm{Fl}(M)$)
\[
 \Phi \mathrel{\tikz[baseline=-.45ex,x=1em,y=1ex]{\draw (0,0)--(2.15,0) (0.32,-0.5)--(0,0) (1.83,0.5)--(2.15,0);}} \Psi \Longleftrightarrow (\Phi\subset\Psi^*)
 \Longleftrightarrow(\Psi\subset\Phi_*) \tag*{(1.1)}
\]
i.e. $\Phi$ a le LLP p.r. à $\Psi$, ou encore $\Psi$ le \deleted{\unsure{R}}RLP
p.r. à $\Phi$. Je \deleted{écris} dis aussi que $(\Phi,\Psi)$ est un
\uline{couple} Q-\uline{orthogonal}. Cette notion n'est pas symétrique,
mais je me permets parfois (exceptionellement) d'écrire
\[
 \Psi \mathrel{\tikz[baseline=-.45ex,x=1em,y=1ex]{\draw (0,0)--(2.15,0) (0.32,0.5)--(0,0) (1.83,-0.5)--(2.15,0);}} \Phi \qquad\text{comme synonyme de}\qquad
 \unsure{\text{``}}\Phi \mathrel{\tikz[baseline=-.45ex,x=1em,y=1ex]{\draw (0,0)--(2.15,0) (0.32,-0.5)--(0,0) (1.83,0.5)--(2.15,0);}} \Psi .
 \tag*{(1.2)}
\]
$\Phi$ est Q-\uline{orthogonal à gauche} à $\Psi$, $\Psi$
\uline{est Q-orthogonal à droite} à $\Phi$, et la Q-orthogonalité
gauche \inserted{(resp. à droite)} est notée par le
\inserted{demi-}sommet des flèches pointant vers le bas
(resp. vers le haut). J'écris
\[
\left\{\begin{aligned}
\Phi \mathrel{\tikz[baseline=-.45ex,x=1em,y=1ex]{\draw (0,0)--(3,0) (0.32,0.5)--(0,0) (0.32,-0.5)--(0,0) (2.68,0.5)--(3,0);}} \Psi &\Longleftrightarrow
 (\Phi\unsure{=|\deleted{=}}\unsure{\Psi^*|\deleted{\Psi^*}}\deleted{\unsure{\Phi|\Psi}}\unsure{)|\deleted{)}}\\
\Phi \mathrel{\tikz[baseline=-.45ex,x=1em,y=1ex]{\draw (0,0)--(3,0) (0.32,-0.5)--(0,0) (2.68,0.5)--(3,0) (2.68,-0.5)--(3,0);}} \Psi &\Longleftrightarrow (\Psi=\Phi_*)\\
\Phi \mathrel{\tikz[baseline=-.45ex,x=1em,y=1ex]{\draw (0,0)--(3,0) (0.32,0.5)--(0,0) (0.32,-0.5)--(0,0) (2.68,0.5)--(3,0) (2.68,-0.5)--(3,0);}} \Psi &\Longleftrightarrow
 (\Phi=\Psi^*\text{ et }\Psi=\Phi_*).
\end{aligned}\right.\tag*{(1.3.)}
\]
Ces trois cas de Q-orthogonalité renforcée sont donc ceux où on a
de plus $\Phi=\Phi^{\sim}$ resp. $\Psi=\overline\Psi$ resp.
$\Phi=\Phi^{\sim}$ et $\Psi=\overline\Psi$. \deleted{\unsure{O}}On dit que $\Phi$ est
Q-\uline{clos à gauche} si $\Phi=\Phi^{\sim}$, que $\Psi$ est
\uline{Q-clos à droite} si $\Psi=\Psi^{-}$, que le couple Q-orthogonal
$(\Phi,\Psi)$ est Q-\uline{clos à gauche} (resp. \uline{à droite})
si $\Phi$ \deleted{\unsure{ou}} (resp. $\Psi$) est Q-clos à g.
(resp. à droite), on dit qu'il est Q-clos si les deux conditions
\scanpage{614}\mspage{2}
% archive sheet number: 613

sont satisfaites simultanément. Quand aucune confusion n'est à
craindre, il m'arrivera de dire \inserted{simplement} que $\Phi$
est Q-clos, ou que $\Psi$ est Q-clos, sans spécifier si c'est à
g. ou à dr. Mais pour un couple Q-orthogonal $(\Phi,\Psi)$ (ou je
préfère l'écrire $\Phi \mathrel{\tikz[baseline=-.45ex,x=1em,y=1ex]{\draw (0,0)--(2.15,0) (0.32,-0.5)--(0,0) (1.83,0.5)--(2.15,0);}} \Psi$), Q-clos signifiera
tjrs Q-clos à gauche et à droite.

Quand un couple Q-orthogonal $\Phi \mathrel{\tikz[baseline=-.45ex,x=1em,y=1ex]{\draw (0,0)--(2.15,0) (0.32,-0.5)--(0,0) (1.83,0.5)--(2.15,0);}} \Psi$ a la
propriété de factorisation de Quillen, je l'indique par la
\deleted{\unsure{ligne}} barre verticale croisant l'horizontale, ainsi
\[
 \Phi \mathrel{\tikz[baseline=-.45ex,x=1em,y=1ex]{\draw (0,0)--(2.6,0) (0.32,-0.5)--(0,0) (2.28,0.5)--(2.6,0) (1.3,-.52)--(1.3,.52);}} \Psi ,\tag*{(1.4)}
\]
et on définit de $=$ le sens des graphismes
\[
\left\{\begin{aligned}
 &\Phi \mathrel{\tikz[baseline=-.45ex,x=1em,y=1ex]{\draw (0,0)--(3.5,0) (0.32,-0.5)--(0,0) (3.18,0.5)--(3.5,0) (1.75,-.52)--(1.75,.52);}} \Psi\\
 &\Phi \unsure{\mathrel{\tikz[baseline=-.45ex,x=1em,y=1ex]{\draw (0,0)--(3.5,0) (0.32,0.5)--(0,0) (0.32,-0.5)--(0,0) (3.18,0.5)--(3.5,0) (3.31,-.25)--(3.5,0) (1.75,-.52)--(1.75,.52);}}|\mathrel{\tikz[baseline=-.45ex,x=1em,y=1ex]{\draw (0,0)--(3.5,0) (0.32,0.5)--(0,0) (0.32,-0.5)--(0,0) (3.18,0.5)--(3.5,0) (1.75,-.52)--(1.75,.52);}}} \Psi\\
 &\Phi \mathrel{\tikz[baseline=-.45ex,x=1em,y=1ex]{\draw (0,0)--(3.5,0) (0.32,-0.5)--(0,0) (3.18,0.5)--(3.5,0) (3.18,-0.5)--(3.5,0) (1.75,-.52)--(1.75,.52);}} \Psi\\
 &\Phi \mathrel{\tikz[baseline=-.45ex,x=1em,y=1ex]{\draw (0,0)--(3.5,0) (0.32,0.5)--(0,0) (0.32,-0.5)--(0,0) (3.18,0.5)--(3.5,0) (3.18,-0.5)--(3.5,0) (1.75,-.52)--(1.75,.52);}} \Psi .
\end{aligned}\right.\tag*{(1.5)}
\]

Un couple Q-\uline{orthogonal} Q-\uline{factorisant}
\deleted{$\Phi,\Psi$} $\Phi \mathrel{\tikz[baseline=-.45ex,x=1em,y=1ex]{\draw (0,0)--(2.6,0) (0.32,-0.5)--(0,0) (2.28,0.5)--(2.6,0) (1.3,-.52)--(1.3,.52);}} \Psi$ est appelé un
\uline{précouple de Quillen}. S'il est de plus Q-clos, i.e.
\deleted{\unsure{lié}} si on a $\Phi \mathrel{\tikz[baseline=-.45ex,x=1em,y=1ex]{\draw (0,0)--(2.6,0) (0.32,0.5)--(0,0) (0.32,-0.5)--(0,0) (2.28,0.5)--(2.6,0) (2.28,-0.5)--(2.6,0) (1.3,-.52)--(1.3,.52);}} \Psi$, on dit que
c'est un \uline{couple de Quillen}.

$\Phi$ est dit \uline{Q-saturé} \uline{à gauche} s'il
\inserted{contient les iso et} est stable par cochangement de base,
par \inserted{(petites)} compositions transfinies, et par facteurs
directs. On \unsure{définit} dualement \deleted{la qualité
\unsure{de} $\Psi$} \deleted{$\unsure{\Phi}$}$\unsure{\Psi|\Phi}$ est dit
\uline{Q-saturé à droite},
\scanpage{615}\mspage{3}
% archive sheet number: 614

contient les iso, stable par chgt de base et par compositions
\unsure{transfinies|transfinie}. Si $\Phi$ est Q-saturé à gauche, il est stable par
\inserted{(petites)} sommes directes i.e. si \inserted{les}
$f_i:A_i\unsure{-|\to}B_i$ ($i\in I$, $I$ petit) sont dans $\Phi$,
$f=\coprod f_i:\coprod A_i\longrightarrow\coprod B_i$ aussi ;
et de $=$ pour \inserted{(petites)} sommes amalgamées finies ou
infinies (quand on se place dans un $A\backslash M$, au lieu de $M$).
On a l'énoncé dual pour $\Psi$, quand $\Psi$ est Q-saturé à droite.
\deleted{\unsure{En fait}}

En fait, il y a des notions plus faibles de Q-saturation que celle
que je viens de dire : Q-\uline{saturation simple} (à gauche :
$\Phi$ contient les iso, stable par composition binaire, par
cochgt de base), Q-\uline{saturation transfinie} (on exige de plus
stabilité par compositions transfinies), enfin Q-\uline{saturation
totale}, ou Q-saturation tout court\unsure{\S|)}, (quand on exige de
plus la stabilité par facteurs directs). On a les implications
\diagram{
\[
\left\{\vcenter{\hbox{
\begin{tikzcd}[ampersand replacement=\&,column sep=0pt,row sep=1em]
\text{Q-}\arrow[d,Rightarrow] \& |[anchor=west]|\text{\deleted{\unsure{sat}}clos à g}\\
\text{Q-}\arrow[d,Rightarrow] \& |[anchor=west]|\text{saturé (total)}\\
\text{Q-}\arrow[d,Rightarrow] \& |[anchor=west]|\text{saturé transfini}\\
\text{Q-} \& |[anchor=west]|\text{saturé simple}\, .
\end{tikzcd}
}}\right.\tag*{(1.6)}
\]
}

\deleted{\unsure{Un}} Pour $\Phi\subset\mathrm{Fl}\,M$, on a sa
Q-\uline{clôture} \inserted{(\uline{à gauche})} $\Phi^{\sim}$
(qui est la plus petite partie Q-close à gauche de
$\mathrm{Fl}\,M$ contenant $\Phi$), et de $=$ son Q-saturé
\inserted{(total)} à gauche, son Q-\uline{saturé transfini}, son
Q-\uline{saturé simple},
\scanpage{616}\mspage{4}
% archive sheet number: 615

avec les \deleted{\unsure{im}}\unsure{implications|inclusions}
tautologiques
\[
\begin{aligned}
\Phi\subset
\underset{\substack{\Vert\\\text{Q-saturé à}\\\text{g. simple}}}{\Phi'}
\subset{}&
\underset{\substack{\Vert\\\text{Q-saturé}\\\text{à g. transfini}}}{\Phi''}
\subset
\underset{\substack{\text{Q-saturé}\\\text{à g. total}}}{\Phi'''}
\subset
\underset{\substack{\text{Q-clôture)}\\\text{à gauche}}}{\Phi^{\sim}}
\end{aligned}\tag*{(1.7)}
\]
et dualement pour $\Psi$. Entre $\unsure{\Phi|\Psi}$ et $\Psi'$
il y a encore un $\Phi'_0$
\[
 \Phi\subset\Phi'_0\subset\Phi',\tag*{(1.8)}
\]
qui est le Q-\uline{saturé pour cochangement de base} (contenant
\inserted{aussi}, par définition, tous les iso). Pour construire
tous ces saturés, on procède ainsi :

1) $\Phi'_0=$ ens des flèches $i$\inserted{$:A\unsure{-|\to}\unsure{D|B}$}
s'insérant dans un carré cocartésien
\diagram{
\begin{tikzcd}[ampersand replacement=\&]
\deleted{\unsure{B}}A_0\arrow[r,"i_0"]\arrow[d] \& B_0\arrow[d]\\
 A\arrow[r,"i"] \& \unsure{D|B}
\end{tikzcd}
},
avec $i_0\in\Phi$ \deleted{\illegible{1}} ou $i_0$ une identité
(auquel cas on trouve pour $i$ les iso).

2) $\Phi'=$ ens des flèches composées multiples
\[
 i_n\,i_{n-1}\,\cdots\,i_1\deleted{\unsure{\xi|\}}}\,,
\]
avec $i_\alpha\in\Phi'_0$ \deleted{\unsure{$\xi$}}.

3) $\Phi''=$ ens des composés transfinis de flèches $\in\Phi'_0$.

\deleted{\unsure{4)}\unsure{$\Phi$}} On a envie de poser
\[
 \begin{aligned}
 \Phi'''=\Phi''\unsure{\natural|{}^{\natural}|\sharp}\overset{\mathrm{def}}{=}
 &\ \text{ens. des facteurs directs}\\
 &\ \text{d'éléments de }\Phi'',
 \end{aligned}
\]
et c'est un fait que $\Phi''\natural\subset\Phi'''$ trivialement, mais
il n'est pas clair en général qu'on ait
\scanpage{617}\mspage{5}
% archive sheet number: 616

égalité, faute de pouvoir montrer que $\Phi''\unsure{\natural|{}^{\natural}|\sharp}$
est stable, non seulement par facteurs directs, mais aussi par cochgt.
de base et par compositions transfinies.

Cependant, rappelons

\uline{\deleted{\unsure{Proposition}}\inserted{Théorème} 1.1}.
Supposons $M$ stable par petites $\varinjlim$, et que ses éléments
soient accessibles. Supposons de plus que $\Phi$ soit petit, plus
généralement qu'il existe une partie $\Phi_0\subset\Phi$ petite telle
que $\Phi\subset\Phi_0^{\sim}$, \deleted{Alors}\inserted{Soit}
\deleted{\unsure{Soit $\Phi$}}
\[
 \Psi\overset{\mathrm{def}}{=}\Phi_*=(\Phi_0)_* .\tag*{(1.1.1)}
\]
\margin{\begin{tabular}{@{}|l@{}}
\uline{NB} On a donc\\
$\Phi \mathrel{\tikz[baseline=-.45ex,x=1em,y=1ex]{\draw (0,0)--(2.15,0) (0.32,-0.5)--(0,0) (1.83,0.5)--(2.15,0) (1.83,-0.5)--(2.15,0);}} \Psi$\\
a fortiori\\
$\Phi_0 \mathrel{\tikz[baseline=-.45ex,x=1em,y=1ex]{\draw (0,0)--(2.15,0) (0.32,-0.5)--(0,0) (1.83,0.5)--(2.15,0) (1.83,-0.5)--(2.15,0);}} \Psi$\end{tabular}}

Sous ces conditions

$\Phi''_0$ désigne le Q-saturé \inserted{à g.} transfini
\deleted{\unsure{de}} de $\Phi_0$, alors on a
$\Phi''_0 \mathrel{\tikz[baseline=-.45ex,x=1em,y=1ex]{\draw (0,0)--(2.6,0) (0.32,-0.5)--(0,0) (2.28,0.5)--(2.6,0) (2.28,-0.5)--(2.6,0) (1.3,-.52)--(1.3,.52);}} \Psi$, i.e. \unsure{toute|\deleted{tout}} flèche $f$ de $M$ se
factorise en $f=pi$, avec $i\in\Phi''_0$, $p\in\Psi$.
% slanted
\margin{Il y a $=$ factorisation \uline{fonctorielle} en $f$,
\deleted{\unsure{pro}} propriété qui mérite une notation, \unsure{spéciale|cf. p. suiv.|cf. p. 64} $\Phi''_0 \mathrel{\tikz[baseline=-.45ex,x=1em,y=1ex]{\draw (0,0)--(2.2,0) (0.32,-0.5)--(0,0) (1.88,0.5)--(2.2,0) (1.88,-0.5)--(2.2,0) (1,-.52)--(1,.52) (1.17,-.52)--(1.17,.52);}} \Psi$}
A fortiori, si $\Phi''$ désigne le Q-saturé \deleted{\unsure{à}} \unsure{à g.|g.} transfini de $\Phi$, on a \deleted{$\unsure{\Phi}$}
\[
 \Phi'' \mathrel{\tikz[baseline=-.45ex,x=1em,y=1ex]{\draw (0,0)--(3,0) (0.32,-0.5)--(0,0) (2.68,0.5)--(3,0) (2.68,-0.5)--(3,0) (1.5,-.52)--(1.5,.52);}} \Psi .
\]

\deleted{\illegible{1}}\uline{Corollaire 1.2} : Sous les conditions
du th, on a
\[
 \Phi'''\ \substack{\text{(\deleted{\unsure{le|Q}}Q-saturé \inserted{à g.} total}\\
                      \text{de }\Phi\text{)}}
 =\underset{\underbrace{\substack{\text{ens des}\\\text{fact. directs}\\
             \text{d'éléments}\\\text{de }\Phi^{\unsure{\prime\prime|o}}}}}
          {\Phi''\natural}
 =\Phi^{\sim} .\tag*{(1.2.1)}
\]

\deleted{\unsure{P}}Plus généralement (en apparence seulement !)

Si $\Phi_1$ est une partie de $\mathrm{Fl}(M)$ qui est Q-saturée à
gauche transfinie, et telle que
\[
 \Phi_0\subset\Phi_1\subset\Phi^{\sim}=\Phi_0^{\sim},
\]
